\documentclass{article}
\usepackage[utf8]{inputenc}
\usepackage{amsmath,amssymb}
\usepackage[most]{tcolorbox}
\usepackage{geometry}
\geometry{margin=2.5cm}

\newcommand{\step}[1]{\par\noindent\hangindent=3.5em\hangafter=1%
  \makebox[3.5em][l]{\textbf{#1.}}}
\newcommand{\steptag}[1]{\hfill{\small\textsf{[#1]}}}

\newtcolorbox{proofbox}[1]{
  colback=gray!5, colframe=gray!60,
  fonttitle=\bfseries, title={#1},
  breakable, enhanced,
  left=4pt, right=4pt, top=4pt, bottom=4pt,
  before skip=10pt, after skip=10pt
}

\begin{document}

\begin{proofbox}{The projected product a•b=P(a∘b) on A satisfies the e...}
\step{1} The projected product a•b=P(a∘b) on A satisfies the easy axioms with O(eta): unit 1•a=a and commutativity a•b=b•a are exact, ||a•b|| <= (1+C eta)||a||||b||, a•a >= -C eta ||a||\textasciicircum{}2 1, and ||a•a|| >= (1-C eta)||a||\textasciicircum{}2. \steptag{validated} \\[6pt]
\step{1.1} EXACT axioms: unit law and commutativity. For a in A: 1•a=P(1 o a). By def-jordan-product 1 o a = (1/2)(1*a+a*1)=a (unit is the algebra identity), so 1•a=P(a)=a since a=P(a) for a in A (def-near-fixed-algebra, A=Im P=Ker(1-P)). Hence 1•a=a EXACTLY. Commutativity: a•b=P(a o b) and b•a=P(b o a); the Jordan product is commutative a o b=b o a (def-jordan-product, (ab+ba)/2 symmetric), so a•b=P(a o b)=P(b o a)=b•a EXACTLY. No eta error in either. \steptag{validated} \\[4pt]
\step{1.2} Upper product bound: ||a•b|| <= (1+C eta)||a|| ||b||. For a,b in A: a•b=P(a o b), so ||a•b||=||P(a o b)|| <= ||P|| * ||a o b|| (definition of operator norm of P, valid since a o b in B(H)\_sa). By GT-Pprops ||P||<=1+C eta. By GT-jb-submult (HOS 3.1.3, V=B(H)\_sa a JB algebra by GT-bhsa-jc) ||a o b|| <= ||a|| ||b||. Combining: ||a•b|| <= (1+C eta) ||a|| ||b||. (C is the universal constant from GT-Pprops, dimension-free.) \steptag{validated} \\[4pt]
\step{1.3} Approx positivity of squares: a•a >= -C eta ||a||\textasciicircum{}2 1. For a in A (self-adjoint, A subset B(H)\_sa), a•a=P(a o a)=P(a\textasciicircum{}2) where a\textasciicircum{}2=a o a. By GT-bh-cstar, a self-adjoint => a\textasciicircum{}2=a\textasciicircum{}*a is positive in B(H), a\textasciicircum{}2>=0. By GT-Pprops delta-positivity (x>=0 => P(x) >= -delta||x|| 1), applied to x=a\textasciicircum{}2>=0: P(a\textasciicircum{}2) >= -delta ||a\textasciicircum{}2|| 1. By GT-jb-axiom-sq (HOS 3.1.4, V=B(H)\_sa JB algebra) ||a\textasciicircum{}2||=||a||\textasciicircum{}2. By GT-Pprops delta<=C eta. Hence a•a=P(a\textasciicircum{}2) >= -delta ||a||\textasciicircum{}2 1 >= -C eta ||a||\textasciicircum{}2 1. \steptag{validated} \\[4pt]
\step{1.4} Square-norm lower bound: ||a•a|| >= (1-C eta)||a||\textasciicircum{}2. For a in A, a•a=P(a\textasciicircum{}2), a\textasciicircum{}2=a o a. STEP A (Jordan-Schwarz): by GT-jordan-schwarz at T=Phi (unital positive on B(H)\_sa), a=a\textasciicircum{}* self-adjoint: Phi(a)\textasciicircum{}2 = Phi(a) o Phi(a) <= Phi(a o a)=Phi(a\textasciicircum{}2). STEP B (norm-monotone in cone): both Phi(a)\textasciicircum{}2>=0 and Phi(a\textasciicircum{}2)>=0 (squares positive, GT-bh-cstar; Phi positive); since 0 <= Phi(a)\textasciicircum{}2 <= Phi(a\textasciicircum{}2), GT-sup-states norm-monotonicity gives ||Phi(a\textasciicircum{}2)|| >= ||Phi(a)\textasciicircum{}2|| = ||Phi(a)||\textasciicircum{}2 (last equality by GT-jb-axiom-sq, ||r\textasciicircum{}2||=||r||\textasciicircum{}2 for r=Phi(a) self-adjoint). STEP C (node 1.4.1): ||Phi(a)|| >= (1-delta)||a||, so ||Phi(a\textasciicircum{}2)|| >= (1-delta)\textasciicircum{}2 ||a||\textasciicircum{}2 >= (1-2 delta)||a||\textasciicircum{}2. STEP D (transfer P<-Phi): ||(P-Phi)(a\textasciicircum{}2)|| <= delta ||a\textasciicircum{}2|| = delta ||a||\textasciicircum{}2 (GT-Pprops operator-norm bound + GT-jb-axiom-sq); by the order-unit norm triangle inequality (GT-orderunit-def), ||P(a\textasciicircum{}2)|| >= ||Phi(a\textasciicircum{}2)|| - ||(P-Phi)(a\textasciicircum{}2)|| >= (1-2 delta)||a||\textasciicircum{}2 - delta||a||\textasciicircum{}2 = (1-3 delta)||a||\textasciicircum{}2 >= (1-C eta)||a||\textasciicircum{}2, using delta<=C eta (GT-Pprops). Hence ||a•a||=||P(a\textasciicircum{}2)|| >= (1-C eta)||a||\textasciicircum{}2. \steptag{validated} \\[6pt]
\step{1.4.1} Reverse-triangle bound (NO ||Phi||=1): ||Phi(a)|| >= (1-delta)||a|| for a in A. Since a in A=Im P, a=P(a) (def-near-fixed-algebra). Then Phi(a)-a = Phi(a)-P(a) = (Phi-P)(a), so ||Phi(a)-a|| = ||(Phi-P)(a)|| <= delta ||a|| where delta:=||P-Phi|| (GT-Pprops: operator-norm bound ||(P-Phi)(z)||<=delta||z||, here z=a; delta<=C eta). By the reverse triangle inequality for the order-unit norm (GT-orderunit-def, the order norm is a genuine norm): ||Phi(a)|| >= ||a|| - ||Phi(a)-a|| >= ||a|| - delta||a|| = (1-delta)||a||. No appeal to ||Phi||=1 or any Phi-contraction. \steptag{validated} \\[4pt]
\end{proofbox}

\end{document}
